\section{Materials and Method}
\label{sec:method}

\subsection{Data}
\label{sec:data}

We use two public parcel-level benchmarks and never tune the protocol between them.

\begin{figure}[tbp]
\centering
\includegraphics[width=\textwidth]{parcelas.pdf}
\caption{Three patches of the held-out fold of the primary benchmark. Top: natural-color
Sentinel-2 composite from the acquisition nearest to 1 July. Bottom: the parcel-level agronomic
annotation. The unit of the product is the parcel, and the catalog this paper studies is the list
of names the bottom row is allowed to use.}
\label{fig:patches}
\end{figure}

\begin{table}[tbp]
\caption{The two benchmarks. Both are described by the same \num{185}-dimensional temporal
feature vector, so the replication varies data, year, region and model family, but not
representation.}
\label{tab:benchmarks}
\centering
\begin{tabular}{lrr}
\toprule
 & primary (PASTIS) & replication (BreizhCrops) \\
\midrule
Season & 2018--2019 & 2017 \\
Region & metropolitan France & Brittany \\
Parcels scored & 16\,640 & 60\,000 \\
Classes & 18 & 9 \\
Majority class & 6\,128 & 18\,113 \\
Minority class & 103 & 2 \\
Imbalance ratio & 60:1 & 9\,057:1 \\
Blocks & 5 spatial (H3, 1\,km buffer) & 2 administrative regions \\
Predictor & pre-trained members & gradient boosting, leave-one-region-out \\
\bottomrule
\end{tabular}
\end{table}

The \textbf{primary benchmark} is PASTIS \cite{garnot2021pastis}. We obtain it through the
PASTIS-R distribution, which adds aligned radar series \cite{saintefaregarnot2022multimodal}, but
\emph{we use its optical modality only}: every predictor here reads Sentinel-2, so calling the
benchmark PASTIS-R would name a sensor we never touch. From it we take the
held-out fold and aggregate the dense annotations to parcel level, yielding \num{16640} parcels
across \num{496} image patches and eighteen agronomic classes. Its imbalance is severe: the
majority class holds \num{6128} parcels and the minority class \num{103}, a ratio of roughly
\num{60}:\num{1}.

The \textbf{replication benchmark} is BreizhCrops \cite{russwurm2020breizhcrops}, Sentinel-2 L2A
series over Brittany for the 2017 season, with nine agronomic classes and a native partition into
administrative regions. We use regions \texttt{frh01} and \texttt{frh04} with a proportional
stratified subsample of \num{30000} parcels each, fixed seed, for a total of \num{60000} parcels.
Its imbalance is an order of magnitude worse than the primary benchmark: in the full index the
majority class outnumbers the minority by up to \num{19207}:\num{1}, and in our subsample three of
the nine classes hold two, seven and \num{289} parcels respectively. Both benchmarks are described
by the same \num{185}-dimensional temporal feature vector, derived from seventeen spectral indices,
so that the two experiments differ in data and model but not in representation.

\begin{figure}[tbp]
\centering
\includegraphics[width=\textwidth]{soporte.pdf}
\caption{Class support on a logarithmic scale. Bars below \num{300} parcels are drawn in the
accent color: these are the classes every retirement criterion removes first. The replication
benchmark's tail is not merely thinner, it is degenerate --- two and seven parcels --- which is
what makes the corridor of Sect.~\ref{sec:corridor} possible.}
\label{fig:support}
\end{figure}

\subsection{Predictors}
\label{sec:predictors}

The protocol operates on posteriors and is agnostic to how they were produced. On the primary
benchmark we take the two strongest of ten available members, a vision transformer for satellite
image time series \cite{tarasiou2023tsvit} and a gradient-boosted model over learned embeddings.
On the replication benchmark we train a gradient-boosted classifier \emph{leaving one region out}:
each parcel's posterior comes from a model that never saw its region, which makes every posterior
we score a genuine out-of-region prediction. That classifier reaches macro-F1 \num{0.4979} and
accuracy \num{0.7421} on \texttt{frh01}, and \num{0.4811} and \num{0.7257} on \texttt{frh04}.

\subsection{Blocks}
\label{sec:blocks}

Every operating point is chosen on data disjoint from the data that measures it. On the primary
benchmark the blocks are spatial: an H3 tessellation at resolution five, clustered and separated by
a \SI{1}{\kilo\meter} buffer, giving five blocks. On the replication benchmark the index publishes
no per-parcel coordinates, so the finest spatial structure available is the administrative region;
we use the two regions as the two blocks. This is a real limitation and we treat it as one: with
two blocks the resampling captures variation \emph{within} regions and not \emph{between} them, so
we report the per-block difference for every contrast rather than a pooled interval alone.

\subsection{The two mechanisms, and a third criterion}
\label{sec:mechanisms}

Let $\mathcal{C}$ be the full set of $C$ classes, $p(\cdot \mid x)$ the predictor's posterior for
parcel $x$, and $\hat{y}(x) = \arg\max_{c \in \mathcal{C}} p(c \mid x)$ its unrestricted
prediction.

\emph{Legend shrinking} selects a subset $\mathcal{L} \subseteq \mathcal{C}$ with
$|\mathcal{L}| = K$ on the training blocks, and on the evaluation block delivers a parcel when
$\hat{y}(x) \in \mathcal{L}$, emitting $\arg\max_{c \in \mathcal{L}} p(c \mid x)$. We test two
criteria for choosing $\mathcal{L}$: by per-class F1 on the training blocks, and by class support
on the training blocks, the latter because it is what practitioners report doing.

\emph{Confidence rejection} keeps $\mathcal{L} = \mathcal{C}$ and delivers the parcels whose
$\max_c p(c \mid x)$ falls above a threshold. The threshold is the quantile of the block's own
confidences that matches the number of parcels the comparison mechanism delivered in that same
block, so the two are compared at equal coverage.

Neither mechanism reads a ground-truth label to decide what to deliver. This is worth stating
explicitly because the natural implementation of legend shrinking does read one: delivering the
parcels whose \emph{true} label lies in $\mathcal{L}$ is an oracle, and it inflates the mechanism's
apparent quality. Our delivery rule uses the unrestricted argmax, which is observable at inference
time.

\subsection{The aligned estimand}
\label{sec:estimand}

Two mechanisms that promise different catalogs cannot be compared by each one's own macro average,
because the averages are taken over different denominators. We therefore score every mechanism with

\begin{equation}
  \mathrm{F1}_{\text{aligned}} \;=\; \frac{1}{|\mathcal{L} \cap \mathcal{Y}|}
  \sum_{c \,\in\, \mathcal{L} \cap \mathcal{Y}} \mathrm{F1}_c ,
  \label{eq:aligned}
\end{equation}

\noindent where $\mathcal{L}$ is the legend under test and $\mathcal{Y}$ the classes actually
present among the delivered parcels. Both mechanisms in a contrast are scored over the \emph{same}
$\mathcal{L}$. We also report each mechanism's native macro over what it individually promises, but
only as a product-facing view; it is not comparable across mechanisms and we never use it for a
contrast.

\subsection{The control}
\label{sec:control}

The control that gives this paper its title is one line. For each legend $\mathcal{L}$ and each
block, score the untouched predictor over $\mathcal{L}$ using Eq.~\ref{eq:aligned} while delivering
\emph{every} parcel. No class is retired from the emission, no parcel is refused; only the
denominator changes. If this control matches or exceeds a mechanism, the quality was bought by the
choice of class set and not by the mechanism.

\subsection{Inference}
\label{sec:inference}

Contrasts use a paired bootstrap: one resample index is drawn per block and \emph{both} mechanisms
are rescored on it, so the two share the sampling noise. We draw \num{1000} resamples and report
the percentile interval of the difference of mean block-level aligned F1, together with the
per-block differences. On the primary benchmark we additionally resample by image patch rather
than by parcel, since parcels within a patch share an image and are not independent. The
pre-registered principal criterion is $K = \lceil C/2 \rceil$; the remaining values of $K$ form an
exploratory family, corrected by the Holm--Bonferroni step-down procedure within each predictor and
each benchmark.
